\chapter{算法、流程与执行路径}

\section{本章作用}

第六章已经把 Pebble 的方法对象与关键规则固定下来。本章继续回答一个更具体的问题：这些方法对象在系统中如何被串联为可执行路径，以及用户从进入页面到得到反馈、保存结果、回流上下文之间，具体经历了哪些状态转换。若说上一章处理的是“系统应如何理解问题”，这一章处理的就是“系统如何把这种理解落实为一条可走通的执行链”。

本章主结论是：Pebble 的执行层由三条主路径与一条横向状态流构成，三条路径分别负责解释、训练与恢复，横向状态流负责关系条件在跨页面场景中的连续注入。

本章关注三条主路径：读心翻译路径、模拟陪练路径与稳态恢复路径。同时，还要补上一条横向的状态流，即“当前关系对象”的选择如何跨页面穿透解码与个人中心，使上下文条件能够随用户移动而保持连续。

\section{总流程概览}

从用户视角看，Pebble 当前的主流程可以压缩为一条入口分流结构：用户先从首页或个人中心进入某一功能面，再根据当前压力状态和任务目标选择不同路径。图~\ref{fig:pebble-workflow-overview} 给出了这一结构的总览。

\begin{figure}[H]
  \centering
  \begin{tikzpicture}[
    node distance=1.55cm and 1.1cm,
    every node/.style={font=\small},
    flowbox/.style={
      rectangle,
      draw=ReportBlueDark,
      rounded corners=8pt,
      fill=ReportSoftBlue,
      minimum width=2.8cm,
      minimum height=1.05cm,
      align=center
    },
    flowalt/.style={
      flowbox,
      fill=ReportSoftGreen
    },
    flowwarn/.style={
      flowbox,
      fill=ReportSoftYellow
    },
    flowarrow/.style={
      -{Latex[length=3mm,width=2mm]},
      thick,
      draw=ReportBlueDark
    }
  ]
    \node[flowbox] (entry) {首页 / 个人中心\\功能入口};
    \node[flowbox, below left=of entry] (translator) {读心翻译\\局部分析};
    \node[flowalt, below=of entry] (dojo) {模拟陪练\\回合反馈};
    \node[flowwarn, below right=of entry] (breathing) {急救呼吸\\稳态恢复};
    \node[flowbox, below=of dojo] (practice) {练习本 / 关系档案\\结果沉淀};

    \draw[flowarrow] (entry) -- (translator);
    \draw[flowarrow] (entry) -- (dojo);
    \draw[flowarrow] (entry) -- (breathing);
    \draw[flowarrow] (translator) -- (practice);
    \draw[flowarrow] (dojo) -- (practice);
    \draw[flowarrow] (breathing) -- (entry);
  \end{tikzpicture}
  \caption{Pebble 的执行路径总览：三条主路径围绕同一入口展开，并在练习沉淀与关系上下文中重新汇流。}
  \label{fig:pebble-workflow-overview}
\end{figure}

这张图表达了三个方法层判断。

\begin{itemize}
  \item 用户入口是统一的。首页负责把功能面并列展示，个人中心负责把关系对象与功能入口重新组织在同一个用户语境下。
  \item 三条主路径分工不同。解码路径强调对当前话语的解释，陪练路径强调回合式练习与即时反馈，呼吸路径强调在高压时刻先恢复状态。
  \item 结果会重新汇流。翻译与陪练的输出都可以进入练习材料，关系对象又会影响下一轮分析，因此整个系统呈现为循环而非单次跳转。
\end{itemize}

\section{路径一：读心翻译的执行链}

（本章使用"读心翻译"作为产品功能名称，其技术实现对应第 3 章定义的"对话解码"过程。）

读心翻译路径是 Pebble 当前最直接的分析入口。其执行顺序可以分为六步。

\begin{itemize}
  \item 用户在翻译页输入需要理解的文本。
  \item 状态层先校验输入是否为空，并把界面状态切换到分析中。
  \item 系统读取当前全局选中的关系对象标识，将其与当前文本一并发送到解码接口。
  \item 解码服务在后端拼接最小必要的关系上下文，并对文本执行结构化分析。
  \item 分析结果返回前端后，被组织为潜台词、情绪状态与三档回应建议。
  \item 用户可直接重试，也可把某一条候选回应存入练习本，作为后续回看材料。
\end{itemize}

这一条路径的关键在于“关系选择”与“解码请求”之间的穿透关系。前端翻译状态并不单独维护一份关系上下文，而是直接读取全局用户中心状态中的当前关系对象标识。这样一来，用户在个人中心完成关系切换后，进入翻译页即可获得上下文增强分析，无需再次重复选择。这一设计让“关系对象作为解释条件”真正进入了解码主路径。

图~\ref{fig:pebble-translate-path-detail} 需要回答的问题是：读心翻译路径在哪些节点发生了状态转换与职责切换。

\begin{figure}[H]
  \centering
  \begin{tikzpicture}[
  node distance=1.15cm,
  every node/.style={font=\small},
  flowstep/.style={
    reportprocess,
    minimum width=6.8cm,
    minimum height=0.85cm,
    align=left,
    text width=6.4cm
  },
  flowdecision/.style={
    reportdecision,
    minimum width=1.4cm,
    minimum height=0.7cm,
    align=center,
    inner sep=1pt
  },
  flowerror/.style={
    reportnode,
    minimum width=2.8cm,
    minimum height=0.75cm,
    align=center
  }
]
  \node[flowstep] (s1) {1. 用户在翻译页输入文本并触发分析};
  \node[flowdecision, below=of s1, xshift=-2.2cm] (d1) {输入\\校验};
  \node[flowerror, below=of s1, xshift=2.2cm] (e1) {返回错误};
  \node[flowstep, below=of d1, xshift=2.2cm] (s2) {2. 读取全局关系标识并注入请求体};
  \node[flowstep, below=of s2] (s3) {3. 调用 POST /api/decode};
  \node[flowdecision, below=of s3, xshift=-2.2cm] (d2) {限流\\校验};
  \node[flowerror, below=of s3, xshift=2.2cm] (e2) {返回 429};
  \node[flowstep, below=of d2, xshift=2.2cm] (s4) {4. 执行脱敏、上下文拼接与分析};
  \node[flowdecision, below=of s4, xshift=-2.2cm] (d3) {LLM\\裁决};
  \node[flowstep, below=of s4, xshift=2.2cm] (s5) {5. 启发式降级};
  \node[flowstep, below=of d3, xshift=2.2cm] (s6) {6. 前端渲染结构化结果};

  \draw[reportarrow] (s1) -- (d1);
  \draw[reportarrow] (d1) -- node[midway, left, font=\tiny] {拒绝} (e1);
  \draw[reportarrow] (d1) -- node[midway, above, font=\tiny] {通过} (s2);
  \draw[reportarrow] (s2) -- (s3);
  \draw[reportarrow] (s3) -- (d2);
  \draw[reportarrow] (d2) -- node[midway, left, font=\tiny] {拒绝} (e2);
  \draw[reportarrow] (d2) -- node[midway, above, font=\tiny] {通过} (s4);
  \draw[reportarrow] (s4) -- (d3);
  \draw[reportarrow] (d3) -- node[midway, above, font=\tiny] {降级} (s5);
  \draw[reportarrow] (d3) -- node[midway, left, font=\tiny] {成功} (s6);
  \draw[reportarrow] (s5) -- (s6);
\end{tikzpicture}

  \caption{读心翻译路径细节图。关键约束：输入校验与限流校验是分析前的连续阻塞决策点；LLM 裁决失败时自动走降级分支。}
  \label{fig:pebble-translate-path-detail}
\end{figure}

读者应带走的判断是：任何阻塞决策点的拒绝都会终止主链，只有全部通过才进入 LLM 分析；降级分支与主分支最终汇流到同一渲染节点，但前端应通过元数据区分二者。

图~\ref{fig:pebble-decode-gate} 回答的问题是：解码路径在哪些节点面临阻塞决策，以及降级路径如何被触发。

\begin{figure}[H]
  \centering
  \begin{tikzpicture}[
  node distance=1.8cm and 1.6cm,
  every node/.style={font=\small}
]
  \node[reportnode] (input) {输入到达};
  \node[reportdecision, right=of input] (validate) {校验与\\限流};
  \node[reportnode, above right=of validate] (llm) {LLM 分析};
  \node[reportnode, below right=of validate] (fallback) {启发式\\降级};
  \node[reportnode, right=of llm] (result) {返回\\DecodeResponse};
  \node[reportnode, below=of result] (error) {返回\\错误响应};

  \draw[reportarrow, ->] (input) -- node[midway, above, font=\tiny] {提交} (validate);
  \draw[reportarrow, ->] (validate) -- node[midway, above left, font=\tiny] {通过} (llm);
  \draw[reportarrow, ->, dashed] (validate) -- node[midway, below left, font=\tiny] {降级} (fallback);
  \draw[reportarrow, ->] (llm) -- node[midway, above, font=\tiny] {成功} (result);
  \draw[reportarrow, ->, dashed] (llm) -- node[midway, right, font=\tiny] {超时/失败} (fallback);
  \draw[reportarrow, ->] (fallback) -- node[midway, below, font=\tiny] {返回} (result);
  \draw[reportarrow, ->] (validate) -- node[midway, left, font=\tiny] {拒绝} (error);
\end{tikzpicture}

  \caption{解码降级门禁。关键约束：LLM 超时或失败时退回到启发式路径；任何降级都必须显式标记在响应元数据中。}
  \label{fig:pebble-decode-gate}
\end{figure}

图~\ref{fig:pebble-decode-gate} 表明解码路径存在两条可返回结果的分支：LLM 分析分支与启发式降级分支。校验与限流裁决属于阻塞点，任何拒绝都会终止后续分析动作，直接返回错误响应。

\section{路径二：模拟陪练的回合流程}

模拟陪练路径承载的是更长的交互链。它相较于读心翻译增加了场景选择、会话状态、轮次消息与即时评分，因此更接近一个轻量的回合式执行系统。

其主流程可以概括为如下步骤。

\begin{itemize}
  \item 页面初始化时加载场景列表，并根据 URL 参数或用户点击结果选定当前场景。
  \item 用户选择场景后，前端创建新会话请求，后端返回 AI 开场白与右侧分析面板的初始信息。
  \item 用户每发送一轮消息，前端先把该消息写入本地消息流，再进入“AI 正在响应”的中间状态。
  \item 后端对这一轮用户输入做情绪与边界分析，同时生成下一轮 AI 回复，并把评分、标签、即时反馈和注意点一起返回。
  \item 前端收到结果后，将 AI 回复附加到消息流，并把右侧面板更新为最新的训练反馈。
  \item 用户可继续发送下一轮，也可选择重启或结束会话。结束后，会话摘要与高频回应可进一步进入练习沉淀路径。
\end{itemize}

这条路径体现出 Pebble 在执行层的一个重要特征：系统不只给出一个静态答案，而是维持一个随轮次更新的状态机。场景、会话、消息流、右侧分析面板与结束摘要共同构成了训练过程中的中间状态，使用户能够在练习中观察自己回应风格的变化。

图~\ref{fig:pebble-dojo-path-detail} 回答的问题是：模拟陪练路径如何把“多轮会话”映射为可持续更新的状态机。

\begin{figure}[H]
  \centering
  \begin{tikzpicture}[
  node distance=1.25cm,
  every node/.style={font=\small},
  flowstep/.style={
    reportprocess,
    minimum width=9.6cm,
    minimum height=0.95cm,
    align=left,
    text width=9.0cm
  }
]
  \node[flowstep] (s1) {1. 用户选择场景并触发会话初始化};
  \node[flowstep, below=of s1] (s2) {2. 前端调用 \texttt{POST /api/simulator}，后端创建会话并返回开场白};
  \node[flowstep, below=of s2] (s3) {3. \texttt{useDojoStore} 写入消息流与右侧分析面板初始状态};
  \node[flowstep, below=of s3] (s4) {4. 用户发送回合消息；后端执行情绪分析并生成下一轮回复};
  \node[flowstep, below=of s4] (s5) {5. 前端更新得分、标签、即时反馈与注意点，进入下一轮循环};
  \node[flowstep, below=of s5] (s6) {6. 结束会话时调用 \texttt{/api/simulator/[sessionId]/end}，返回摘要并可沉淀到练习本};

  \draw[reportarrow] (s1) -- (s2);
  \draw[reportarrow] (s2) -- (s3);
  \draw[reportarrow] (s3) -- (s4);
  \draw[reportarrow] (s4) -- (s5);
  \draw[reportarrow] (s5) -- (s6);
\end{tikzpicture}

  \caption{模拟陪练路径细节图，突出会话启动、回合更新与结束摘要的状态切换。}
  \label{fig:pebble-dojo-path-detail}
\end{figure}

该图的结论是：陪练路径的核心在于跨回合反馈对象的连续更新能力。

\section{路径三：稳态恢复的短流程}

急救呼吸路径承担的是与前两条路径不同的职责。它不以语言分析为中心，而以状态恢复为中心，因此执行流程更短、交互目标更明确。

\begin{itemize}
  \item 用户进入呼吸页后，系统立即开始倒计时与呼吸节奏动画。
  \item 页面围绕 4-7-8 周期自动切换吸气、屏息与呼气阶段，并用视觉律动同步提示当前阶段。
  \item 用户可随时暂停、继续或在完成后重新开始。
  \item 当用户状态恢复后，再回到首页、翻译页或陪练页进入后续理解与应对路径。
\end{itemize}

这一条路径的意义在于，它为整个系统提供了一条“先恢复，再分析”的执行分支。对于高压状态下的用户，立即给出语言分析并不总是最优入口；先稳定状态，再回到分析与练习链路，更符合第五章中“先理解后应对”与“边界克制”的设计要求。

图~\ref{fig:pebble-breathing-path-detail} 回答的问题是：稳态恢复路径为何可在本地状态机内闭环完成。

\begin{figure}[H]
  \centering
  \begin{tikzpicture}[
  node distance=1.25cm,
  every node/.style={font=\small},
  flowstep/.style={
    reportprocess,
    minimum width=9.6cm,
    minimum height=0.95cm,
    align=left,
    text width=9.0cm
  }
]
  \node[flowstep] (s1) {1. 用户进入急救呼吸页，界面启动倒计时与节奏提示};
  \node[flowstep, below=of s1] (s2) {2. \texttt{useBreathing} 调用呼吸状态机，进入吸气-屏息-呼气循环};
  \node[flowstep, below=of s2] (s3) {3. 动画层与文案层按阶段同步更新，提供实时节奏反馈};
  \node[flowstep, below=of s3] (s4) {4. 用户可暂停、继续或重置，控制恢复过程的交互节奏};
  \node[flowstep, below=of s4] (s5) {5. 倒计时结束后返回首页、翻译页或陪练页继续主任务};

  \draw[reportarrow] (s1) -- (s2);
  \draw[reportarrow] (s2) -- (s3);
  \draw[reportarrow] (s3) -- (s4);
  \draw[reportarrow] (s4) -- (s5);
\end{tikzpicture}

  \caption{稳态恢复路径细节图，突出“状态机循环”和“可控暂停重置”机制。}
  \label{fig:pebble-breathing-path-detail}
\end{figure}

该图的结论是：呼吸路径在工程上可以独立于后端资源运行，从而保证高压时刻的最小可用性。

\section{跨模块时序：解码链如何穿越前后端边界}

图~\ref{fig:pebble-workflow-cross-sequence} 需要回答的问题是：同一轮解码请求在跨模块传播时遵循怎样的调用顺序。

\begin{figure}[H]
  \centering
  \begin{tikzpicture}[font=\small]
  \node[reportnode, minimum width=2.0cm, minimum height=0.85cm] at (0, 0) {用户};
  \node[reportnode, minimum width=2.0cm, minimum height=0.85cm] at (2.6, 0) {翻译页};
  \node[reportnode, minimum width=2.0cm, minimum height=0.85cm] at (5.2, 0) {TranslatorStore};
  \node[reportdecision, minimum width=1.6cm, minimum height=0.7cm] at (7.8, 0) {校验\\门禁};
  \node[reportnode, minimum width=2.0cm, minimum height=0.85cm] at (10.4, 0) {DecodeService};

  \foreach \x in {0,2.6,5.2,7.8,10.4} {
    \draw[densely dashed, draw=ReportGray] (\x, -0.45) -- (\x, -8.0);
  }

  \draw[reportarrow] (0, -1.0) -- node[above, font=\tiny]{1. 输入文本} (2.6, -1.0);
  \draw[reportarrow] (2.6, -1.7) -- node[above, font=\tiny]{2. 触发解码} (5.2, -1.7);
  \draw[reportarrow] (5.2, -2.4) -- node[above, font=\tiny]{3. 读取关系标识} (5.2, -2.4);
  \draw[reportarrow] (5.2, -3.1) -- node[above, font=\tiny]{4. POST /api/decode} (7.8, -3.1);
  \draw[reportarrow] (7.8, -3.8) -- node[above, font=\tiny]{5. 通过} (10.4, -3.8);
  \draw[-{Latex[length=3mm,width=2mm]}, thick, dashed, draw=ReportBlueDark]
    (7.8, -4.5) -- node[above, font=\tiny]{拒绝 400/429} (5.2, -4.5);
  \draw[reportarrow] (10.4, -5.2) -- node[above, font=\tiny]{6. 调用 analyzeText} (10.4, -5.2);
  \draw[-{Latex[length=3mm,width=2mm]}, thick, dashed, draw=ReportBlueDark]
    (10.4, -5.9) -- node[above, font=\tiny]{7. DecodeResponse} (7.8, -5.9);
  \draw[-{Latex[length=3mm,width=2mm]}, thick, dashed, draw=ReportBlueDark]
    (7.8, -6.6) -- node[above, font=\tiny]{8. 响应体} (5.2, -6.6);
  \draw[-{Latex[length=3mm,width=2mm]}, thick, dashed, draw=ReportBlueDark]
    (5.2, -7.3) -- node[above, font=\tiny]{9. 更新状态并渲染} (2.6, -7.3);
\end{tikzpicture}

  \caption{解码链跨模块时序图。关键约束：校验门禁是 LLM 分析前的阻塞决策点；任何拒绝裁决都会终止后续分析动作，降级裁决则把请求导入启发式路径。}
  \label{fig:pebble-workflow-cross-sequence}
\end{figure}

读者应带走的判断是：校验门禁属于主路径同步裁决。若校验返回拒绝，前端 Store 应直接进入错误状态并提示用户；若校验通过但 LLM 超时，系统自动走降级路径并显式标记降级元数据，前端据此区分高质量分析与 fallback 结果。

\section{横向状态流：关系对象如何跨路径保持连续}

除了三条主路径之外，Pebble 还有一条横向的状态流，即当前关系对象的跨页面保持。个人中心负责加载和展示当前关系，关系选择结果会被持久化保存；翻译路径在发起解码请求时直接读取这一状态；个人中心再次打开时又会根据保存的标识重新加载关系详情。

这一横向状态流带来三点执行层收益。

\begin{itemize}
  \item 用户无需在每个页面重复指定同一关系对象，交互成本更低。
  \item 关系信息进入了解码主路径，分析结果能够更稳定地体现情境差异。
  \item 个人中心、关系图谱与核心功能面之间形成统一的用户语境，系统不再表现为若干互不相干的工具页。
\end{itemize}

\section{执行边界与失败处理}

Pebble 当前的执行路径还包含一组必要的边界处理。它们决定系统在异常、空输入、会话失效或用户未完成前置动作时如何退回安全状态。

\begin{itemize}
  \item 解码路径在空输入时直接返回错误状态，避免无意义请求进入分析服务。
  \item 陪练路径要求会话已启动且消息非空，否则会阻断发送并给出错误提示。
  \item 会话结束后，陪练路径不再接受继续发送，而是要求用户显式重启或重新开始。
  \item 呼吸路径允许暂停与重启，从而把恢复过程设计成可控节奏，而非强制性流程。
\end{itemize}

这些边界处理虽然并不直接构成产品卖点，却是执行路径成立的前提。没有这些状态约束，系统很容易在多轮交互中出现上下文漂移、空请求分析或会话错位，从而削弱整个闭环的稳定性。

\section{本章小结}

本节读完后，读者应能画出读心翻译、模拟陪练与急救呼吸三条主路径的执行顺序，说明关系对象如何在跨页面时保持连续注入，并指出解码路径中至少两个阻塞决策点。
